\section{Related Work}
\label{sec:related}

\para{Alignment and diversity loss.}
The observation that alignment reduces output diversity has been documented across
multiple settings. \citet{kirk2023rlhf} measured per-input and across-input mode
collapse after RLHF on LLaMA-7B, finding substantial reductions in distinct n-gram
diversity and Sentence-BERT similarity. Counterintuitively, increasing the KL
penalty in PPO did not recover diversity---it reduced it further while also
degrading performance. \citet{lin2023mitigating} characterized the broader
``alignment tax'' on downstream benchmarks and found that simple weight-space
interpolation between pre- and post-RLHF checkpoints ($\theta_{\text{merged}} =
(1-\interpcoeff) \cdot \theta_{\text{SFT}} + \interpcoeff \cdot
\theta_{\text{RLHF}}$) outperformed all other mitigation strategies, including
KL penalties, LoRA, and experience replay. Our work differs from these studies
in two ways: we focus specifically on \emph{stylistic} diversity (targeted
impersonation of distinct voices) rather than generic output diversity, and we
operate in logit space rather than weight space.

\para{KL-constrained alignment.}
DPO~\cite{rafailov2023dpo} implicitly optimizes a KL-constrained reward objective
with the closed-form solution $\pi^*(y|x) \propto \pi_{\text{ref}}(y|x) \cdot
\exp(r(x,y)/\beta)$, where $\beta$ controls divergence from the reference policy.
\citet{korbak2022rl} showed that this framework is equivalent to Bayesian
inference: the base model serves as the prior, the reward provides the likelihood,
and the aligned model is the posterior. Larger $\beta$ preserves more of the prior
distribution. Extensions to alternative f-divergences~\cite{wang2023fdpo,
go2023fdivergence} demonstrated that Jensen-Shannon divergence often provides
better alignment--diversity tradeoffs than reverse KL. We build on this
understanding by asking: regardless of which divergence was used during training,
can we navigate between base and aligned distributions \emph{post hoc}?

\para{Distribution arithmetic and interpolation.}
Several lines of work have explored combining base and aligned model distributions
at inference time. \citet{zhou2024emulated} introduced emulated disalignment,
showing that alignment can be arithmetically inverted via the log-linear
combination $\pi_{\text{ED}}(y_t) \propto \pi_{\text{base}}(y_t)^{\interpcoeff+1}
/ \pi_{\text{align}}(y_t)^{\interpcoeff}$, effectively doubling the harmfulness
of base models. Contrastive decoding~\cite{li2023contrastive} and proxy
tuning~\cite{mitchell2023proxy} similarly combine model distributions to steer
generation. Most recently, \citet{zhu2025adding} added a single
identity-initialized layer trained with DPO, enabling continuous interpolation
along a $\lambda$ axis at inference time with only 2\% of parameters trained.
\citet{lu2024online} integrated model merging into every RLHF training step,
showing that online weight averaging improves both preference scores and benchmark
performance. Our contribution applies distribution arithmetic specifically to
stylistic capabilities, testing whether these distributions are continuously
navigable.

\para{Superficial Alignment Hypothesis.}
\citet{zhou2023lima} proposed that alignment primarily teaches format and tone
rather than new knowledge, suggesting that the base model's capabilities persist
largely intact. \citet{ji2024resistance} provided mechanistic evidence: aligned
models exhibit ``elasticity,'' reverting toward the pretraining distribution upon
further fine-tuning at a rate proportional to the ratio of pretraining to alignment
data. \citet{lake2024distributional} further showed that alignment shifts diversity
from distributional (diverse across samples) to Overton (diverse within a single
response), and that base model behavior is recoverable via in-context learning.
Our results provide direct evidence for the Superficial Alignment Hypothesis in the
domain of stylistic generation: base distributions survive alignment, but alignment
adds a crucial \emph{addressing mechanism} for accessing them controllably.
